%======================================================================
\section{Coefficients of the evolution operator up to $\order{g^{2}}$ through LCPT}
\label{sec:coeff}
%======================================================================

To determine the coefficients of the evolution operator we compare its action on the elementary
Fock states, Eqs.~\eqref{eq:Omvac}--\eqref{eq:Omqg}, with the light-cone wave functions computed in
section~\ref{sec:omega}, using the rule \eqref{eq:rule}. Since these states probe different
sectors of Fock space, they constrain different structures of $\OM$, and the coefficients are fixed
one by one. Wherever possible we express the second order coefficients through the first order
ones by the reduction formula \eqref{eq:reduction}, which in terms of coefficients reads: for a
transition $i\to f$ through an intermediate state $m$ generated by the first order structures with
coefficients $X$ ($i\to m$) and $Y$ ($m\to f$),
\begin{equation}
K_{f\leftarrow i}\supset Y\,X\;\frac{E_{m}-E_{f}}{E_{i}-E_{f}},
\qquad\text{and}\qquad
K_{f\leftarrow i}\supset\frac{(2\pi)^{3(n+M)/2}\,H^{(2)}_{fi}}{E_{i}-E_{f}}
\label{eq:redcoef}
\end{equation}
for the instantaneous terms, the internal momenta being integrated with $\int_{p}$. This is the
form in which the coefficients $D_{1}$, $D_{2}$ of Paper~I were written. In the terms that involve the
three-gluon coefficients $C$ and $C_{2}=-C^{\dg}$, the combinatorial factors associated with identical
gluons are those of Paper~I, whose conventions for these coefficients we adopt throughout.

\subsection{Order $g$}

Matching \eqref{eq:Pg}, \eqref{eq:Pq-qg} and \eqref{eq:Pqb-qbg} with \eqref{eq:Omg} and \eqref{eq:Omq}
($K=(2\pi)^{9/2}\psi$) gives
\begin{align}
\mf A^{b\alpha\beta}_{3\,j\lambda_{2}\lambda_{1}}(p,m,l)&=
\frac{(2\pi)^{3}g\,t^{b}_{\beta\alpha}}{2\sqrt2\sqrt{p^{+}}}\;
\frac{\dthree(p-l-m)\;\Ph^{j}_{\lambda_{1}\lambda_{2}}(p;l,m)}{\eps_{p}-\eps_{l}-\eps_{m}},
\label{eq:A3}\\
\mf F^{\beta'\beta b}_{\lambda'\lambda j}(l',l,p)&=
\frac{(2\pi)^{3}g\,t^{b}_{\beta\beta'}}{2\sqrt2\sqrt{p^{+}}}\;
\frac{\dthree(l'-l-p)\;\Ph^{j}_{\lambda\lambda'}(p;l,l')}{\eps_{l'}-\eps_{l}-\eps_{p}},
\label{eq:F}\\
\mf G^{\alpha'\alpha b}_{\lambda'\lambda j}(m',m,p)&=
-\frac{(2\pi)^{3}g\,t^{b}_{\alpha'\alpha}}{2\sqrt2\sqrt{p^{+}}}\;
\frac{\dthree(m'-m-p)\;\Ph^{j}_{\lambda'\lambda}(p;m',m)}{\eps_{m'}-\eps_{m}-\eps_{p}} .
\label{eq:G}
\end{align}
Each of them is the vertex \eqref{eq:V3}--\eqref{eq:VG} divided by the energy released in the
splitting, the quark analogue of the three-gluon coefficient
$C=(2\pi)^{9/2}\bra{gg}H_{ggg}\ket{g}/(E_{i}-E_{f})$ of Paper~I. Their conjugates
$\mf A_{4}$, $\mf F_{2}$, $\mf G_{2}$ follow from \eqref{eq:U1}; for $\mf A_{4}$ this was verified directly
in \eqref{eq:Pqq-g}.

\subsection{Pair creation off the valence charge: $\mf A_{1}$}

Comparing \eqref{eq:Omvac} with \eqref{eq:Pvac} ($K=(2\pi)^{3}\psi$),
\begin{equation}
\boxed{\;
\mf A^{\alpha\beta}_{1\,\lambda_{2}\lambda_{1}}(m,l)=
-g^{2}\,t^{a}_{\beta\alpha}\,\rho^{a}(-k)\left[
\frac{\delta_{\lambda_{1}\lambda_{2}}}{k^{+}\bm k^{2}}
+\frac{\chi^{\dg}_{\lambda_{1}}\big(\bm l\!\cdot\!\bm m+i\,(\bm l\!\times\!\bm m)\,\sigma^{3}\big)\chi_{\lambda_{2}}}
{2\,l^{+}m^{+}\bm k^{2}\,(\eps_{l}+\eps_{m})}\right],\qquad k=l+m,\;}
\label{eq:A1}
\end{equation}
with $\bm l\!\times\!\bm m\equiv\epsilon^{ij}l^{i}m^{j}$. In the form \eqref{eq:redcoef},
$\mf A_{1}=\int_{k}\mf A_{3}(k,m,l)A(k)\,(\eps_{l}+\eps_{m}-\eps_{k})/(\eps_{l}+\eps_{m})
-(2\pi)^{3}\bra{q\bar q}H_{q-{\rm inst}}\vac/(\eps_{l}+\eps_{m})$, with the
Weizs\"acker--Williams coefficient of Paper~I
\begin{equation}
A^{b}_{j}(p)=-\frac{\sqrt2\,g\,p^{j}\rho^{b}(-p)}{\sqrt{p^{+}}\,\bm p^{2}} .
\label{eq:Acoef}
\end{equation}
\emph{Check.} With \eqref{eq:A1}, $\mf A_{1}+\mf A_{2}^{\dg}$ obtained from the direct LCPT
computation of section~\ref{sec:wf-qq} equals $\int_{p}\mf A_{3}(p,m,l)A(p)$, which is the
relation \eqref{eq:UA2}. Using $4(k^{+})^{2}\eps_{k}=2k^{+}\bm k^{2}$, this amounts to the identity
$\frac{1}{2k^{+}\bm k^{2}E}+\frac{1}{4(k^{+})^{2}E(E-\eps_{k})}=\frac{1}{2k^{+}\bm k^{2}(E-\eps_{k})}$
with $E=\eps_{l}+\eps_{m}$; the instantaneous parts cancel.

\subsection{Pair creation off a soft gluon: $\mf B_{3}$}

Comparing \eqref{eq:Omg} with the connected wave functions of section~\ref{sec:wf-g}
($K=(2\pi)^{6}\psi$),
\begin{equation}
\boxed{\;
\mf B^{\alpha\beta ad}_{3\,\lambda_{2}\lambda_{1}in}(m,l,k,r)=A^{d}_{n}(r)\,\mf A^{a\alpha\beta}_{3\,i\lambda_{2}\lambda_{1}}(k,m,l)
+(2\pi)^{6}\Big[\psi^{(1)}+\psi^{(2)}+\psi^{(3)}+\psi^{(6)}+\psi^{(7)}\Big]_{gq\bar q},\;}
\label{eq:B3}
\end{equation}
where $\psi^{(n)}_{gq\bar q}$ denote the coefficient functions of the $q\bar q g$ components
\eqref{eq:B3a}, \eqref{eq:B3b}, \eqref{eq:B3c}, \eqref{eq:B3f} and \eqref{eq:B3g}, i.e.\ the expressions
multiplying $\ket{q\bar q\,g}$ under the integral; their prefactors become $(2\pi)^{3}g^{2}/8$ for
(a), (b), (c) and $(2\pi)^{3}g^{2}/2$ for (f), (g). The first term is the factorized contribution
\eqref{eq:B3de}. Equivalently, by \eqref{eq:redcoef},
\begin{align}
\mf B_{3}(m,l,k,r)=\;&A(r)\mf A_{3}(k,m,l)+\frac{1}{D_{f}}\Big[\int_{p}\mf F(p,l,r)\mf A_{3}(k,m,p)\,(\eps_{p}-\eps_{l}-\eps_{r})
\nonumber\\
&+\int_{m'}\mf G(m',m,r)\mf A_{3}(k,m',l)\,(\eps_{m'}-\eps_{m}-\eps_{r})
\nonumber\\
&+\int_{k'}\mf A_{3}(k',m,l)\,C(k,k',r)\,(\eps_{k'}-\eps_{l}-\eps_{m})
\nonumber\\
&+(2\pi)^{6}\bra{q\bar q\,g}H_{ggq\bar q}+H_{q-{\rm inst}}\ket{g}\Big],
\label{eq:B3red}
\end{align}
with $D_{f}=\eps_{k}-\eps_{l}-\eps_{m}-\eps_{r}$. In each product the weight cancels the energy
denominator of the coefficient that acts second, so that \eqref{eq:B3red} is (vertex)$\times$(first
order coefficient)$/D_{f}$, as in \eqref{eq:B3a}--\eqref{eq:B3c}.

\subsection{Gluon pair annihilation into quarks: $\mf B_{1}$}

The two-gluon state probes the symmetric combination $\mf B_{1}(p,k)+\mf B_{1}(k,p)$,
Eq.~\eqref{eq:Omgg}. Assigning to $\mf B^{ba\alpha\beta}_{1\,ji\lambda_{2}\lambda_{1}}(p,k,m,l)$ the
contributions in which the gluon $\gl{a}{i}{k}$ is absorbed (by the valence charge, by the quark or
by the antiquark) while $\gl{b}{j}{p}$ splits, and half of the contributions that are symmetric by
themselves,
\begin{equation}
\boxed{\;
\begin{aligned}
\mf B^{ba\alpha\beta}_{1\,ji\lambda_{2}\lambda_{1}}(p,k,m,l)=\;&
-A^{\dg a}_{i}(k)\,\mf A^{b\alpha\beta}_{3\,j\lambda_{2}\lambda_{1}}(p,m,l)
\\&+(2\pi)^{6}\Big[\tfrac12\psi^{(1)}+\psi^{(4)}+\psi^{(5)}+\tfrac12\psi^{(6)}+\tfrac12\psi^{(7)}\Big]_{q\bar q},
\end{aligned}\;}
\label{eq:B1}
\end{equation}
where $\psi^{(1)},\psi^{(6)},\psi^{(7)}$ are the coefficient functions of \eqref{eq:B1a}, \eqref{eq:B1e},
\eqref{eq:B1f}, and $\psi^{(4)},\psi^{(5)}$ those of \eqref{eq:B1c}, \eqref{eq:B1d} with
$(a,i,k)\leftrightarrow(b,j,p)$ in the displayed term. The first term is the factorized
contribution \eqref{eq:B1b}. In the form \eqref{eq:redcoef} the $t$- and $u$-channel terms read
\begin{multline}
\mf B_{1}(p,k,m,l)\supset\frac{1}{D_{f}}\Big[\int_{l_{1}}\mf F_{2}(k,l_{1},l)\,\mf A_{3}(p,m,l_{1})\,(\eps_{l_{1}}+\eps_{k}-\eps_{l})
\\+\int_{m_{1}}\mf G_{2}(k,m_{1},m)\,\mf A_{3}(p,m_{1},l)\,(\eps_{m_{1}}+\eps_{k}-\eps_{m})\Big],
\label{eq:B1red}
\end{multline}
which makes manifest that they produce the $\mf F$ and $\mf G$ terms of the unitarity relation
\eqref{eq:UB2}. Since $D_{f}=\eps_{k}+\eps_{p}-\eps_{l}-\eps_{m}$ can vanish, the $\rho$ independent part
of $\mf B_{1}$ has an on-shell pole whose residue is the $gg\to q\bar q$ amplitude.

\subsection{Factorized coefficients: $\mf C_{3}$, $\mf C_{4}$, $\mf D_{1}$}

For these sectors the transition proceeds through two independent subprocesses, the two time
orderings combine by \eqref{eq:factor}, and the coefficients factorize. With
$K=(2\pi)^{9}\psi$ and $(2\pi)^{9}/64\pi^{3}=(2\pi)^{6}/8$, Eqs.~\eqref{eq:Pggqq}, \eqref{eq:Pgqq}
and \eqref{eq:Pqqqq} give
\begin{align}
\mf C^{bcde\alpha\beta}_{4\,jkon\lambda_{2}\lambda_{1}}(p,q,r,s,m,l)&=
\mf A^{b\alpha\beta}_{3\,j\lambda_{2}\lambda_{1}}(p,m,l)\;C^{edc}_{nok}(q,r,s),
\label{eq:C4}\\
\mf C^{bcde\alpha\beta}_{3\,jkon\lambda_{2}\lambda_{1}}(p,q,r,s,m,l)&=
\mf A^{b\alpha\beta}_{3\,j\lambda_{2}\lambda_{1}}(p,m,l)\;C^{edc}_{2\,nok}(q,r,s),
\label{eq:C3}\\
\mf D^{cb\alpha'\beta'\alpha\beta}_{1\,kj\lambda_{2}'\lambda_{1}'\lambda_{2}\lambda_{1}}(q,p,m',l',m,l)&=
\mf A^{b\alpha\beta}_{3\,j\lambda_{2}\lambda_{1}}(p,m,l)\;\mf A^{c\alpha'\beta'}_{3\,k\lambda_{2}'\lambda_{1}'}(q,m',l'),
\label{eq:D1}
\end{align}
where $C$ and $C_{2}=-C^{\dg}$ are the $g\to gg$ and $gg\to g$ coefficients of Paper~I
(Eqs.~(4.8) and (5.5) there). Explicitly, e.g.,
\begin{equation}
\mf C_{4}=\frac{i(2\pi)^{6}g^{2}}{8}\,t^{b}_{\beta\alpha}f^{cde}\,
\frac{\dthree(p-l-m)\dthree(q-r-s)}{\sqrt{p^{+}q^{+}r^{+}s^{+}}}\;
\frac{\Ph^{j}_{\lambda_{1}\lambda_{2}}(p;l,m)\,\Vg^{kon}(q;r,s)}
{\big(\eps_{p}-\eps_{l}-\eps_{m}\big)\big(\eps_{q}-\eps_{r}-\eps_{s}\big)} .
\end{equation}
The factorizations hold with no leftover powers of $2\pi$, which is a stringent check of the
prefactors of $\mf A_{3}$ and of the rule \eqref{eq:rule}. The unitarity partners
$\mf C_{1}$, $\mf C_{2}$, $\mf D_{2}$ follow from \eqref{eq:UC1}--\eqref{eq:UD2}; since $\mf A_{3}$ contains no
$\rho$, two $\mf A_{3}$'s commute and \eqref{eq:UD2} gives
$\mf D_{2}(l,m,l',m',q,p)=\mf A^{\dg}_{3}(q,m,l)\mf A^{\dg}_{3}(p,m',l')$, the crossed assignment, which is
also what a direct LCPT computation of $q\bar q\,q\bar q\to gg$ gives.

\subsection{Quark--antiquark scattering: $\mf D_{3}$}

Comparing \eqref{eq:Omqq} with \eqref{eq:D3a}--\eqref{eq:D3e} ($K=(2\pi)^{6}\psi$; with the ordering
$d\,b$ of \eqref{eq:Om2qq} no extra sign arises),
{\small
\begin{equation}
\boxed{\;
\begin{aligned}
&\mf D^{\beta'\alpha'\alpha\beta}_{3\,\lambda_{1}'\lambda_{2}'\lambda_{2}\lambda_{1}}(l',m',m,l)
=\frac{(2\pi)^{3}g^{2}}{8}\,\frac{\dthree(l+m-l'-m')}{D_{f}}\Bigg\{
\\&\quad
t^{a}_{\beta\alpha}t^{a}_{\alpha'\beta'}\,
\frac{\Ph^{i}_{\lambda_{1}\lambda_{2}}(k;l,m)\,\big[\Ph^{i}_{\lambda_{1}'\lambda_{2}'}(k;l',m')\big]^{*}}
{k^{+}\big(\eps_{l'}+\eps_{m'}-\eps_{k}\big)}
\\
&\quad-t^{d}_{\beta\beta'}t^{d}_{\alpha'\alpha}\Bigg[
\frac{\Ph^{n}_{\lambda_{1}\lambda_{1}'}(r;l,l')\,\Ph^{n}_{\lambda_{2}'\lambda_{2}}(r;m',m)\,\theta(r^{+})}
{r^{+}\big(\eps_{l'}-\eps_{l}-\eps_{r}\big)}
\\&\hspace{2.6cm}+\frac{\Ph^{n}_{\lambda_{1}\lambda_{1}'}(\tilde r;l,l')\,\Ph^{n}_{\lambda_{2}'\lambda_{2}}(\tilde r;m',m)\,\theta(\tilde r^{+})}
{\tilde r^{+}\big(\eps_{m'}-\eps_{m}-\eps_{\tilde r}\big)}\Bigg]
\\
&\quad-\frac{8\,t^{a}_{\beta\beta'}t^{a}_{\alpha'\alpha}\,\delta_{\lambda_{1}\lambda_{1}'}\delta_{\lambda_{2}\lambda_{2}'}}{(l'^{+}-l^{+})^{2}}
+\frac{8\,t^{a}_{\beta\alpha}t^{a}_{\alpha'\beta'}\,\delta_{\lambda_{1}\lambda_{2}}\delta_{\lambda_{1}'\lambda_{2}'}}{(k^{+})^{2}}\Bigg\},
\end{aligned}}
\label{eq:D3}
\end{equation}}
with $k=l+m$, $r=l'-l$, $\tilde r=m'-m$ and $D_{f}=\eps_{l'}+\eps_{m'}-\eps_{l}-\eps_{m}$. In the form
\eqref{eq:redcoef},
\begin{align}
\mf D_{3}(l',m',m,l)=\frac{1}{D_{f}}\Big[&\int_{k}\mf A_{3}(k,m,l)\mf A_{4}(l',m',k)\,(\eps_{k}-\eps_{l}-\eps_{m})
\nonumber\\&+\int_{r}\mf G_{2}(r,m',m)\mf F(l',l,r)\,(\eps_{m'}+\eps_{r}-\eps_{m})
\nonumber\\
&+\int_{\tilde r}\mf F_{2}(\tilde r,l',l)\mf G(m',m,\tilde r)\,(\eps_{l'}+\eps_{\tilde r}-\eps_{l})
\nonumber\\&+(2\pi)^{6}\bra{q\bar q}H_{q-{\rm inst}}\ket{q\bar q}\Big].
\label{eq:D3red}
\end{align}
\emph{Check.} Write $D_{a}=\eps_{k}-\eps_{l}-\eps_{m}$ and $D_{a}'=\eps_{k}-\eps_{l'}-\eps_{m'}$, so that
$D_{a}-D_{a}'=D_{f}$. The $s$-channel term of $\mf D_{3}(l',m',m,l)+\mf D_{3}^{\dg}(l,m,m',l')$ is
\begin{multline}
\frac{(2\pi)^{3}g^{2}}{8}\,\frac{t^{a}t^{a}\,\Ph\Ph^{*}\dthree}{k^{+}}\Big[-\frac{1}{D_{a}'D_{f}}+\frac{1}{D_{a}D_{f}}\Big]
\\=-\frac{(2\pi)^{3}g^{2}}{8}\,\frac{t^{a}t^{a}\,\Ph\Ph^{*}\dthree}{k^{+}D_{a}D_{a}'}
=-\int_{k}\mf A_{3}(k,m,l)\,\mf A_{3}^{\dg}(k,m',l'),
\end{multline}
which is the first term of \eqref{eq:UD3}. In the same way the $t$-channel terms (b) and (c)
reproduce the $\mf F$, $\mf G$ terms of \eqref{eq:UD3}, while the instantaneous term is anti-Hermitian
and drops out. Note that the $t$-channel part of $\mf D_{3}$ is not anti-Hermitian by itself; it
satisfies unitarity only together with the coefficients $\mf F$ and $\mf G$.

\subsection{Quark rescattering on the valence charge: $\mf P$ and $\bar{\mf P}$}

Comparing \eqref{eq:Omq} with \eqref{eq:Pq-q} ($K=(2\pi)^{3}\psi$, $(2\pi)^{3}/32\pi^{3}=1/4$),
\begin{equation}
\boxed{\;
\begin{aligned}
\mf P^{\beta'\beta}_{\lambda'\lambda}(l',l)=\frac{g^{2}}{4}\,\frac{t^{b}_{\beta\beta'}\,\rho^{b}(l'-l)}{(p^{+})^{2}}\,
\chi^{\dg}_{\lambda}\Bigg\{&\theta(l'^{+}-l^{+})\left[\frac{2}{\Delta-\eps_{p}}
-\frac{p^{+}\,(\bm\sigma\!\cdot\!\bm l)(\bm\sigma\!\cdot\!\bm l')}{l^{+}l'^{+}(\Delta-\eps_{p})\,\Delta}\right]
\\
&-\theta(l^{+}-l'^{+})\left[\frac{2}{\eps_{p}}
-\frac{p^{+}\,(\bm\sigma\!\cdot\!\bm l)(\bm\sigma\!\cdot\!\bm l')}{l^{+}l'^{+}\,\eps_{p}\,\Delta}\right]\Bigg\}\chi_{\lambda'},
\end{aligned}}
\label{eq:P}
\end{equation}
with $\Delta=\eps_{l'}-\eps_{l}$ and $p=\pm(l'-l)$ in the two regions. In the form \eqref{eq:redcoef},
\begin{multline}
\mf P(l',l)=-A^{\dg}(p)\,\mf F(l',l,p)\,\frac{\eps_{p}}{\Delta}\Big|_{p=l'-l}
\\-\mf F^{\dg}(l,l',p)\,A(p)\,\frac{\eps_{l'}+\eps_{p}-\eps_{l}}{\Delta}\Big|_{p=l-l'}
+\frac{(2\pi)^{3}\bra{q(l)}H_{q-{\rm inst}}\ket{q(l')}}{\Delta} .
\label{eq:Pred}
\end{multline}
The antiquark coefficient $\bar{\mf P}^{\alpha'\alpha}_{\lambda'\lambda}(m',m)$ is $(2\pi)^{3}$ times the
coefficient function of \eqref{eq:Pqb-qb}.

\emph{Check.} Let $p=l'-l$ with $l'^{+}>l^{+}$. Then $\mf P(l',l)$ is given by the first line of
\eqref{eq:P}, and $\mf P(l,l')$ (incoming $l$, outgoing $l'$) by the second line with $l\leftrightarrow l'$.
Using $\rho^{b}(-p)^{\dg}=\rho^{b}(p)$ and $\big(t^{b}_{\beta'\beta}\big)^{*}=t^{b}_{\beta\beta'}$, and denoting
$S=\chi^{\dg}_{\lambda}(\bm\sigma\!\cdot\!\bm l)(\bm\sigma\!\cdot\!\bm l')\chi_{\lambda'}/(l^{+}l'^{+})$,
\begin{align}
\mf P(l',l)+\mf P^{\dg}(l,l')&=\frac{g^{2}}{4}\frac{t^{b}\rho^{b}(p)}{(p^{+})^{2}}
\Big\{2\delta_{\lambda\lambda'}\Big[\frac{1}{\Delta-\eps_{p}}-\frac{1}{\eps_{p}}\Big]
-p^{+}S\Big[\frac{1}{(\Delta-\eps_{p})\Delta}+\frac{1}{\eps_{p}\Delta}\Big]\Big\}
\nonumber\\
&=\frac{g^{2}}{4}\frac{t^{b}\rho^{b}(p)}{(p^{+})^{2}}\,
\frac{4\eps_{p}-2\Delta-p^{+}S}{(\Delta-\eps_{p})\,\eps_{p}}
\nonumber\\
&=\frac{g^{2}}{4}\frac{t^{b}\rho^{b}(p)\;p^{j}\Ph^{j}_{\lambda\lambda'}(p;l,l')}{(p^{+})^{2}(\Delta-\eps_{p})\,\eps_{p}}
=-\int_{p'}A^{\dg}(p')\,\mf F(l',l,p'),
\end{align}
where we used \eqref{eq:pGam1}, \eqref{eq:Acoef} and \eqref{eq:F}. This is the relation \eqref{eq:UP} (the
self-energy term is momentum diagonal and scaleless). The pole at $\Delta=0$ cancels in the sum, as
it must: $\mf P+\mf P^{\dg}$ is a product of first order coefficients.

\subsection{Coefficients with three outgoing particles: $\mf Q_{1}$, $\mf Q_{2}$, $\mf K_{1}$}

Comparing \eqref{eq:Omq}, \eqref{eq:Omqg} with \eqref{eq:Pq-qgg}, \eqref{eq:Pq-qqq} and with the
quark--gluon amplitude of section~\ref{sec:wf-qg}, and using \eqref{eq:redcoef}, we obtain
\begin{align}
\mf Q_{1}(l',l,q,r)=\frac12\Bigg\{&A(q)\mf F(l',l,r)+A(r)\mf F(l',l,q)
\nonumber\\&+\frac{1}{D_{f}}\Big[\int_{l_{1}}\mf F(l_{1},l,q)\mf F(l',l_{1},r)\,(\eps_{l_{1}}-\eps_{l}-\eps_{q})+(q\leftrightarrow r)
\nonumber\\
&+\int_{s}C(s,q,r)\,\mf F(l',l,s)\,(\eps_{s}-\eps_{q}-\eps_{r})
\nonumber\\&+(2\pi)^{6}\bra{q\,gg}H_{ggq\bar q}+H_{q-{\rm inst}}\ket{q}\Big]\Bigg\},
\label{eq:Q1}
\end{align}
with $D_{f}=\eps_{l'}-\eps_{l}-\eps_{q}-\eps_{r}$;
\begin{align}
\mf K_{1}(l',l,l_{2},m)=\frac12\Bigg\{&\frac{1}{D_{f}}\Big[\int_{k}\mf A_{3}(k,m,l_{2})\,\mf F(l',l,k)\,(\eps_{k}-\eps_{l_{2}}-\eps_{m})
\nonumber\\&+(2\pi)^{6}\bra{q\,q\bar q}H_{q-{\rm inst}}\ket{q}\Big]-\big(l\leftrightarrow l_{2}\big)\Bigg\},
\label{eq:K1}
\end{align}
with $D_{f}=\eps_{l'}-\eps_{l}-\eps_{l_{2}}-\eps_{m}$; and
\begin{align}
\mf Q_{2}(l',q,l,r)=&-A^{\dg}(q)\mf F(l',l,r)+A(r)\mf F_{2}(q,l',l)
\nonumber\\
&+\frac{1}{D_{f}}\Big[\int_{l_{1}}\mf F(l_{1},l,r)\mf F_{2}(q,l',l_{1})\,(\eps_{l_{1}}-\eps_{l}-\eps_{r})
\nonumber\\&\qquad+\int_{l_{1}}\mf F_{2}(q,l_{1},l)\mf F(l',l_{1},r)\,(\eps_{l_{1}}+\eps_{q}-\eps_{l})
\nonumber\\
&\qquad-\int_{m_{1}}\mf A_{4}(l',m_{1},r)\mf A_{3}(q,m_{1},l)\,(\eps_{l'}+\eps_{m_{1}}-\eps_{r})
\nonumber\\&\qquad+\int_{s}\mf F_{2}(s,l',l)\,C(q,s,r)\,(\eps_{s}+\eps_{l'}-\eps_{l})
\nonumber\\
&\qquad+\int_{s}C_{2}(s,q,r)\,\mf F(l',l,s)\,(\eps_{s}+\eps_{q}-\eps_{r})
\nonumber\\&\qquad+(2\pi)^{6}\bra{q\,g}H_{ggq\bar q}+H_{q-{\rm inst}}\ket{q\,g}\Big],
\label{eq:Q2}
\end{align}
with $D_{f}=\eps_{l'}+\eps_{q}-\eps_{l}-\eps_{r}$. The six terms in the bracket of \eqref{eq:Q2} are the
$s$-channel, the two orderings of the $u$-channel (the second, with a soft pair in the intermediate
state, carries the Fermi minus sign) and the two orderings of the $t$-channel, followed by the
instantaneous terms; the first line contains the factorized $\rho$ dependent products. The
instantaneous matrix elements are given in appendix~\ref{app:me}. The antiquark coefficients
$\bar{\mf Q}_{1,2}$, $\bar{\mf K}_{1}$ follow by $\mf F\to\mf G$, $\mf A_{3}\to\mf A_{3}$ with the roles of
quark and antiquark exchanged.

\subsection{Coefficients fixed by unitarity}

The remaining coefficients follow from the constraint equations of section~\ref{sec:omega}:
$\mf A_{4}$, $\mf F_{2}$, $\mf G_{2}$ from \eqref{eq:U1}; $\mf A_{2}$ from \eqref{eq:UA2}; $\mf B_{2}$, $\mf B_{4}$
from \eqref{eq:UB2}, \eqref{eq:UB4}; $\mf C_{1}$, $\mf C_{2}$, $\mf D_{2}$ from \eqref{eq:UC1}--\eqref{eq:UD2};
$\mf Q_{3}$ and $\mf K_{2}$ from \eqref{eq:UQ3}, \eqref{eq:UK2}. Finally, every structure of
$\OM^{(2)}_{\times}$, Eq.~\eqref{eq:Omx}, is the normal ordered product of two first order structures
$X$, $Y$ acting on disjoint particles; by \eqref{eq:factor}, and since at most one of the two
coefficients contains $\rho$,
\begin{equation}
K_{:\hat X\hat Y:}=X\,Y,
\label{eq:Kx}
\end{equation}
e.g.\ the coefficient of $a^{\dg}b^{\dg}d^{\dg}b^{\dg}b\,a$ is $\mf A_{3}\,\mf F$ and that of
$a^{\dg}a^{\dg}b^{\dg}b\,a^{\dg}a$-type structures is $\mf F\,C$. These coefficients satisfy the unitarity
relations identically.

The overall energy denominator $E_{i}-E_{f}$ can vanish only in the coefficients that preserve
the number of soft quarks, antiquarks and gluons, namely $\mf P$, $\bar{\mf P}$, the $\rho$ independent
parts of $\mf D_{3}$ and $\mf Q_{2}$, and in the $2\to2$ coefficients $\mf B_{1}$, $\mf B_{2}$. These are
understood with the prescription of section~\ref{sec:lcpt}, and their residues are on-shell
scattering amplitudes. In all other coefficients the overall denominator either never vanishes
(away from collinear configurations of measure zero) or cancels between the two time orderings,
which then combine into a product of first order coefficients, as in $\mf C_{3}$, $\mf C_{4}$,
$\mf D_{1}$, the factorized parts of $\mf B_{1}$, $\mf B_{3}$, $\mf Q_{1}$, $\mf Q_{2}$, and \eqref{eq:Kx}.

\subsection{The generator}
\label{sec:generator}

Having determined all the coefficients, we reconstruct the evolution operator. As in Paper~I we
write $\OM=e^{iG}$ with $G=G_{\rm YM}+G_{q}$ Hermitian and known through $\order{g^{2}}$. Since
$\OM^{(2)}=iG^{(2)}+\tfrac12(iG^{(1)})^{2}$ and the uncontracted part of $\tfrac12(iG^{(1)})^{2}$ is
$\OM^{(2)}_{\times}$ together with $\mf C_{1,\dots,4}$, $\mf D_{1,2}$ and the factorized parts of the other
coefficients, the connected part of $G_{q}$ is
\begin{align}
G_{q}=&-i\int_{p,l,m}\Big[\mf A_{3}\,b^{\dg}d^{\dg}a-\mf A_{3}^{\dg}\,a^{\dg}db\Big]
-i\int_{p,l,l'}\Big[\mf F\,a^{\dg}b^{\dg}b-\mf F^{\dg}\,b^{\dg}b\,a\Big]
-i\int_{p,m,m'}\Big[\mf G\,a^{\dg}d^{\dg}d-\mf G^{\dg}\,d^{\dg}d\,a\Big]
\nonumber\\
&-\frac{i}{2}\int_{l,m}\Big[\mf A_{1}-\mf A_{2}^{\dg}\Big]b^{\dg}d^{\dg}
-\frac{i}{2}\int_{l,m}\Big[\mf A_{2}-\mf A_{1}^{\dg}\Big]db
-\frac{i}{2}\int_{p,q,l,m}\Big[\mf B_{1}-\mf B_{2}^{\dg}\Big]b^{\dg}d^{\dg}aa
\nonumber\\
&-\frac{i}{2}\int_{p,q,l,m}\Big[\mf B_{2}-\mf B_{1}^{\dg}\Big]a^{\dg}a^{\dg}db
-\frac{i}{2}\int_{p,q,l,m}\Big[\mf B_{3}-\mf B_{4}^{\dg}\Big]a^{\dg}a\,b^{\dg}d^{\dg}
\nonumber\\
&-\frac{i}{2}\int_{p,q,l,m}\Big[\mf B_{4}-\mf B_{3}^{\dg}\Big]db\,a^{\dg}a
\nonumber\\
&-\frac{i}{2}\int_{l,m,l',m'}\Big[\mf D_{3}-\mf D_{3}^{\dg}\Big]b^{\dg}d^{\dg}db
-\frac{i}{2}\int_{l,l'}\Big[\mf P-\mf P^{\dg}\Big]b^{\dg}b
-\frac{i}{2}\int_{m,m'}\Big[\bar{\mf P}-\bar{\mf P}^{\dg}\Big]d^{\dg}d
\nonumber\\
&-\frac{i}{2}\int\Big[\mf Q_{1}-\mf Q_{3}^{\dg}\Big]a^{\dg}a^{\dg}b^{\dg}b
-\frac{i}{2}\int\Big[\mf Q_{3}-\mf Q_{1}^{\dg}\Big]b^{\dg}b\,aa
-\frac{i}{2}\int\Big[\mf Q_{2}-\mf Q_{2}^{\dg}\Big]a^{\dg}b^{\dg}b\,a
\nonumber\\
&-\frac{i}{2}\int\Big[\mf K_{1}-\mf K_{2}^{\dg}\Big]b^{\dg}b^{\dg}d^{\dg}b
-\frac{i}{2}\int\Big[\mf K_{2}-\mf K_{1}^{\dg}\Big]b^{\dg}d\,bb
\nonumber\\
&+\big(q\to\bar q\big),
\label{eq:Gq}
\end{align}
where the arguments and labels are those of \eqref{eq:Om1q}--\eqref{eq:OmF}, $(q\to\bar q)$ denotes the
antiquark analogues of the $\mf Q$ and $\mf K$ terms, and $\mf F$, $\mf G$, $\mf A_{3}$ are of order $g$ while
all the other coefficients are of order $g^{2}$. $G_{q}$ is Hermitian by construction. By the
unitarity relations, each combination $[K-K'^{\dg}]/2$ in \eqref{eq:Gq} equals $K$ minus half of the
corresponding product of first order coefficients; for the particle number preserving coefficients
it is their anti-Hermitian part, which contains the instantaneous terms and the principal value
parts of the on-shell poles.
